\section{Twisted reciprocity and actual trace products}\label{sec:reciprocity}
The opposite class has coefficient action $\chi^{-1}$. Its ray relations therefore use the inverse cyclotomic eigenspace, not the ordinary character rows that will appear later for differences of primitives.

In the descended all inert regime, let $M=\Q(\zeta_7)$ and let $\gamma\in H^1(\Q,\F_7(\chi^{-1}))$ restrict to the relevant beta class. Over $M$, the restriction is additive and satisfies
\[
\gamma(\delta_a g\delta_a^{-1})=a^{-1}\gamma(g).
\]
At a conductor prime $q\equiv1\pmod7$, there are six degree one primes in $M$. At $q\equiv-1\pmod7$, there are three degree two primes. If $z_q$ is a chosen seventh root of unity and $\ell_q$ is the corresponding seventh power logarithm, the local functional at the embedding $\zeta_7\mapsto z_q^b$ is $b s_q\ell_q$.

For degree two, the indices $b$ and $-b$ represent the same prime. Frobenius sends a residue $x$ to $x^q$, and $\ell_q(x^q)=-\ell_q(x)$. Therefore
\[
(-b)\ell_q(x^q)=b\ell_q(x).
\]
The functional is well defined on prime ideals, not just on the six embeddings. Principal reciprocity gives, for a test prime generator $\pi_l$ outside the conductor,
\begin{equation}\label{eq:orbitreciprocity}
R_{l,q}=-\sum_{b\in\mathrm{reps}_q}b\ell_q(\pi_l(z_q^b))
=-f_q^{-1}\sum_{b=1}^6b\ell_q(\pi_l(z_q^b)),
\end{equation}
where $f_q=1$ or two. Using all six terms without dividing by two in the second case would double count each prime. Constant weights or inverse $b$ weights give different modules.

\subsection{The same representation trace product}
At depth one and an unramified place of residue norm $Q\equiv1\pmod7$, write the actual Frobenius in the fixed natural frame as
\[
F=\begin{pmatrix}a&b\\7c&d\end{pmatrix},\qquad\det F=Q.
\]
Put $d=1+7D$ and $\tr(F)=Q+1+7T$. Then $a=Q+7(T-D)$, and the determinant equation gives
\begin{equation}\label{eq:trace-product-exact}
bc=T+(Q-1)D+7TD-7D^2.
\end{equation}
Modulo seven this becomes
\begin{equation}\label{eq:trace-product}
\frac{\tr(F)-Q-1}{7}=\alpha(F)\beta(F)\quad\text{in }k.
\end{equation}
This is the coefficient field rank two trace of the actual representation. It is not half the rational four dimensional trace. The hypothesis $Q\equiv1\pmod7$ makes the extra diagonal term disappear. At other residue norms, \eqref{eq:trace-product-exact} must be retained.

At good places, the right side must equal one of the corrections extracted from the common quartic \eqref{eq:quartic}. At bad toric places, exact trace $Q+1$ makes the product zero. The zero product condition does not always imply $\beta(F)=0$, because $\alpha(F)$ may already vanish. This is the error corrected in the Block 70 draft. It was consequential: imposing the stronger zero beta condition would falsely remove valid necessary system data.

If the actual polynomial is
\[
P(X)=X^4-SX^3+BX^2-QSX+Q^2,
\]
the correction $u=(\tau-Q-1)/7$ has
\[
\operatorname{Tr}(u)=\frac{S-2(Q+1)}7,\qquad
\operatorname{Nm}(u)=\frac{(B-2Q)-(Q+1)S+(Q+1)^2}{49}
\]
modulo seven, under the stated divisibility assumptions. The two implementations extract these values using different methods, including an integer trace discriminant and an enumeration of $k$ by trace and norm. Neither assigns a Frobenius value at a beta conductor prime by pretending beta is unramified there. Such conductor values require the separate self unit calculation of Section~\ref{sec:self-conductor}.

\subsection{Coefficient rigidity without radical identification}
The recovered coefficient theorem, initially in the descended depth one regime and later in the all depth semilinear formulation, normalizes the natural class as
\begin{equation}\label{eq:rigidalpha}
\alpha=vw[\phi]+r[3]+[\sqrt5],\qquad v,r\in\F_7^\times.
\end{equation}
Its mechanism retains the actual joint finite quotient. The three Kummer classes are independent after cyclotomic restriction. Distinct cyclotomic eigenspaces separate the natural and opposite images. Chebotarev then realizes joint vectors at good split primes, where the same quartic at the paired $K$ places forces products to be coefficient conjugate. The resulting conditions are $u^7=-u$ and $r^7=r$. In the descended subcase beta has an $\F_7$ representative.

This is a cohomological coefficient statement. It does not identify the base radical with $w$ and does not descend an arbitrary Tate lattice. In the more general semilinear formulation, beta may decompose as a descended plus component and a minus component twisted by the quadratic character of $K/\Q$. The bounded all inert computation uses the stipulated plus regime. The minus component remains a separate arithmetic obligation.
